\Titular*%
{historia}%
{\textit{Pikadon:} cando o Sol baixa á Terra}%
{Santiago González Gómez}%
{Que lle ocorre a un humano na zona cero dunha explosión nuclear?
Terroríficos debuxos dos superviventes de Hiroshima reflicten o inferno vivido en agosto de 1945.}%

\noindent\fbox{
    \parbox{0.95\linewidth}{
        \textbf{AVISO}: Os feitos aquí relatados e as imaxes incluídas poden resultar altamente
        perturbadores ao lector. É necesario coñecer os horrores da guerra nuclear para
        evitala nun futuro, pero aconséllase á xente sensible ler con precaución.
    }
}

\begin{multicols}{2}

O 6 de agosto de 1945, Os Estados Unidos soltaron a primeira bomba atómica
sobre Hiroshima, O Xapón. Tres días despois, unha segunda bomba caeu en Nagasaki.
Estes bombardeos nucleares, que saldaron con centos de miles de mortes,
desataron unha destrución nunca antes vista na Terra e precipitaron a rendición
do Imperio Xaponés, dando fin á Segunda Guerra Mundial.

Foron as bombas un mal menor? Eran necesarias? Nunca o saberemos. Porén, o
horror que experimentaron os habitantes de Hiroshima e Nagasaki, especialmente
aqueles que non morreron inmediatamente, é innegable. Imos tentar presentar
aquí que lle ocorreu á xente que se atopaba en Hiroshima aquel fatídico día.

\begin{figure}[H]
    \centering
    \includegraphics[width=\linewidth]{revistas/003/imaxes/fungo.jpg}
    \caption{
        Fungo nuclear visto dende as aforas de Hiroshima. Susumu Horikoshi, 6
        anos no bombardeo.
    }
    \label{fig:fungo}
\end{figure}

\subsection*{O inferno na terra}

\textit{Little Boy}, a bomba que estourou sobre Hiroshima, non detonou ao
chegar o chan, senón 600 m sobre a superficie. Unha detonación aérea asegura
que a onda expansiva chegue o máis lonxe posible. \textit{Little Boy} estaba
armada con 64 kg de uranio-235. O funcionamento da bomba, a grandes liñas, era
o seguinte: mediante o uso de explosivos convencionais, un pequeno proxectil
deste uranio era lanzado contra un maior obxectivo cilíndrico de uranio. Ao
xuntárense, acadábase a masa crítica, iniciando a reacción nuclear. A explosión
resultante foi equivalente a uns 15 quilotóns de TNT.

Os habitantes de Hiroshima puideron ver un cegador resplandor branco (de aquí
vén o alcume \textit{pikadon}, literalmente luz estrondosa, que os xaponeses
deron ás bombas), seguido dunha enorme onda expansiva. A presión xerada por
\textit{Little Boy} non era suficiente para matar unha persoa, pero si para
lanzala varios metros polos aires e provocarlle importantes feridas. Tamén era
suficiente para derrubar todos os edificios nun raio dun quilómetro e medio.
Agás certas estruturas de formigón, Hiroshima foi completamente aplanada. Gran
parte da xente morreu esmagada polos entullos ou quedou irremediablemente
atrapada.

Parte da enerxía da explosión, liberada en forma calorífica, xerou unha bóla de
lume cunha temperatura superficial similar á do Sol, incendiando todo material
inflamable preto da zona cero. A temperatura era suficiente para descolorar
permanentemente a rocha, de tal xeito que persoas, postes, bicicletas... que
tapaban certas partes dos edificios deixaron sombras permanentes impresas.

\begin{figure}[H]
    \centering
    \includegraphics[width=1\linewidth]{revistas/003/imaxes/carbon.jpg}
    \caption{
        Corpos carbonizados. Unha nai cun bebé mantense ergueita, correndo, na
        posición que tiña ao morrer. Yasuko Yamagata, 17 anos no bombardeo.
    }
    \label{fig:carbon}
\end{figure}

Con respecto ás persoas que se atopaban suficientemente preto da zona cero,
todas morreron de xeito inmediato. Aínda que aqueles que se atopaban moi, moi
preto, si foron vaporizados (que posiblemente é a imaxe que temos moitos de nós
do resultado dunha bomba atómica), a meirande parte foron carbonizados; é
dicir, foron expostos a unha calor tal (e durante un tempo suficiente) que os
seus tecidos se deshidrataron e foron convertidos en carbón.

\begin{figure}[H]
    \centering
    \includegraphics[width=\linewidth]{revistas/003/imaxes/procesion.jpg}
    \caption{
        Procesión de vítimas queimadas coa pel colgando. Kichisuke Yoshimura,
        18 anos no bombardeo.
    }
    \label{fig:procesion}
\end{figure}

Isto tamén lle sucedeu á xente que quedou atrapada e sucumbiu aos numerosos
lumes provocados pola bomba. Aquelas persoas que se atopaban lixeiramente máis
lonxe sufriron queimaduras profundas, provocadas pola calor arrastrada pola
onda expansiva dende o epicentro. A algúns derretéronlles a carne e os globos
oculares. Estas queimaduras foron capaces de incinerar a maior parte da roupa;
aquela que non desapareceu deixou en moitas ocasións os seus padróns gravados
para sempre na pel de debaixo.

Os que sobreviviron ás súas queimaduras e conseguiron saír dos entullos
atopáronse camiñando entre a desolación. A mestura da intensa calor e a onda
expansiva foi capaz de arrincar a pel dos brazos, pernas ou costas de moitos
dos superviventes, deixando expostos músculos e ósos. As queimaduras foron
especialmente graves na cara e nos brazos, que eran en moitas ocasións as
partes máis expostas. De xeito instintivo, para evitar que lles quedasen
pegados, os superviventes comezaron a camiñar cos seus brazos estendidos,
similarmente á imaxe que temos os occidentais dos zombis; e os xaponeses, das
pantasmas.

En resumo, as rúas de Hiroshima enchéronse de procesións de mortos en vida,
brazos adiante, coa pel colgando da punta dos dedos ou arrastrándose detrás
deles. Estas ``Santas Compañas'' de xente dirixíanse cara aos ríos, xa que a
calor era insoportable e á maioría asaltoulles unha sede inhumana. As
testemuñas describen á xente botándose aos ríos, que baixaban cheos de
cadáveres. Moitos tamén beberon; a maioría destes morreron pola contaminación
da auga.

\begin{figure}[H]
    \centering
    \includegraphics[width=\linewidth]{revistas/003/imaxes/rio.jpg}
    \caption{
        Cruzando o río cheo de cadáveres. Toshiko Kihara, 17 anos no bombardeo.
    }
    \label{fig:rio}
\end{figure}

\subsection*{A radiación}

Como a detonación foi aérea, non houbo cinza radioactiva (\textit{fallout}) na
zona cero. Os restos de radioisótopos da explosión ascenderon á atmosfera ata
albergárense na estratosfera. Este é o proceso durante o cal se forma a famosa
nube con forma de fungo. Entre media hora e unha hora despois da explosión,
estes restos condensaron o vapor de auga e comezaron a precipitar, de volta á
cidade, en forma dunha chuvia negra, espesa e lodosa. De novo, moitos dos
afectados beberon da choiva, morrendo polos efectos da radiación.

\begin{figure}[H]
    \centering
    \includegraphics[width=\linewidth]{revistas/003/imaxes/cremacion.jpg}
    \caption{
        Nai busca un lugar no que cremar o seu fillo calcinado. Matsumuro
        Kazuo, 32 anos no bombardeo.
    }
    \label{fig:cremacion}
\end{figure}

O feito de que non se depositaren grandes cantidades de restos radiactivos na
cidade é o que provocou que a xente puidese repoboar o terreo de xeito case
inmediato, a diferenza de Pripiat (Chernóbil), que segue a ser inhabitable a
día de hoxe. Porén, os presentes no momento do bombardeo si foron irradiados
con partículas alfa, beta e gamma. Moitos dos que recibiron unha dose letal de
radiación morreron antes polas queimaduras ou os entullos que por esta; porén,
aqueles que conseguiran sobrevivir ao impacto inicial morreron semanas despois
con cadros dolorosísimos que inclúen caída do pelo, vómitos e escumación negra
na boca, febres altas, diarrea e hemorraxias xeneralizadas. Os médicos non
comprendían esta morte súbita e silenciosa que sobreviña a xente que parecera
ilesa pola explosión.

\begin{figure}[H]
    \centering
    \includegraphics[width=\linewidth]{revistas/003/imaxes/choiva.jpg}
    \caption{
        Muller sedenta bebe da choiva negra. Takakura Akio, 19 anos no
        bombardeo.
    }
    \label{fig:choiva}
\end{figure}

Moitos dos que recibiron doses radiactivas non mortais acabaron desenvolvendo
enfermidades graves a longo prazo, especialmente cancros (en particular
leucemia). As forzas de ocupación estadounidenses estudaron durante anos os
efectos radioactivos en moitas destas persoas, que se converteron en cobaias
humanas. Porén, nunca comunicaban aos afectados as enfermidades que estaban
desenvolvendo, e moito menos proporcionaron algún tipo de tratamento. Os
estadounidenses tamén prohibiron á prensa falar sobre os bombardeos e os seus
efectos. Cando as forzas de ocupación marcharon, a desinformación entre o resto
da poboación xaponesa era tal que os superviventes (\textit{hibakusha}) foron
amplamente discriminados, por medo a que a súa simple presenza puidese
contaxialos con algunha enfermidade. Moitos superviventes tamén optaron por non
ter fillos, temendo as posibles malformacións ou enfermidades que poderían
causarlles.

\subsection*{Debuxos dende a zona cero}

Moitas das cousas relatadas polos \textit{hibakusha} poden parecer
inverosímiles. As cámaras non chegaron ata Hiroshima ata días despois, e moito
material foi censurado logo polos estadounidenses. Con todo, algúns
superviventes animáronse a debuxar as imaxes gravadas na súa mente, que nos
permiten facernos unha idea de que viviron aquel día. Todas as imaxes están
sacadas de \cite{HPMMP:Abomb}, e constitúen só unha pequena porción, a menos
perturbadora, de todos os debuxos producidos polos \textit{hibakusha}. Anímase
o lector a ver máis imaxes das gardadas en \cite{HPMMP:Abomb}.

\Cita[Tadayoshi Saika, \\ No cenotafio de Hiroshima]{ Descansade en paz, pois non repetiremos o erro. }

Loitemos a prol do desarme nuclear.

\nocite{MIT:Ground}
\nocite{BBC:Atomic}

\printbibliography

\end{multicols}
